% =====================================================================
% CHAPTER 5: Experimental Setup and Data Collection (new department
% template, 2026-07-14). Red sections fixed: Introduction (first),
% Summary (last). Black middle sections are ours; the template suggests
% Experimental Setup, Test Bench, Instruments, Calibration, Experimental
% Procedure, and Data Collection, which this chapter follows closely.
% Assembled per thesis/RESTRUCTURE_PLAN_V2.md (ch5 owner) from:
%   - old mainmatter/ch7.tex Section "Simulation and Experimental Setup"
%     (7.6; fig:setup, fig:jetson intact) and Section "Testing and
%     Measurement" (7.8; \label{sec:testing_protocols} kept).
%   - old mainmatter/ch2.tex "Benchmarks and Datasets" subsection (frozen
%     dataset numbers reused; the table itself, tab:datasets, stays in
%     ch2 as literature survey material, so a new table with a distinct
%     label is built here for the datasets this thesis actually runs on).
%   - old mainmatter/ch3.tex Principle-of-Operation paragraph on the rig's
%     calibrated focal length and baseline (numbers only; the disparity-
%     to-depth mechanism itself stays with ch:implementation).
%   - old appendices/app_a.tex "Software and Hardware Environment" and
%     "Reproduction Pointers" (compute/software facts; script names).
% New prose: 5.1 Introduction, 5.7 Summary. No new numbers introduced;
% every value here is copied from the sources above with its citation.
% =====================================================================
\chapter{Experimental Setup and Data Collection}
\label{ch:experimental}

\section{Introduction}
Chapter~\ref{ch:implementation} described how the StereoLite network was
built, trained, and optimized. Before its numbers can be trusted, the
conditions under which they were produced need to be stated precisely:
what hardware ran the training and the timing runs, what data the
network saw during training and what it never saw, where the real
camera pairs came from and how the rig behind them was calibrated, and
what procedure turned a forward pass into a reported number. This
chapter fixes that record. It covers the compute platforms and the
physical stereo rig that together form the test bench, the calibration
that makes the rig's geometry usable for depth recovery, the three data
sources the thesis draws on, how the real-camera pairs were collected
and how a frozen teacher network supplied a comparison reference for
them, and the four measurement protocols that produce every accuracy,
latency, and robustness figure this thesis reports.

The chapter is organized to move from apparatus to data to procedure.
Section~\ref{sec:exp_setup_bench} lists the compute platforms and the
stereo rig that make up the test bench. Section~\ref{sec:calibration}
gives the rig's calibrated parameters and the rectification convention
they support. Section~\ref{sec:exp_datasets} describes the three data
sources in turn: the synthetic training and test corpus, the zero-shot
real benchmark, and the physical rig's own captures. Section
\ref{sec:data_collection} details how the real pairs were captured and
how pseudo-ground-truth was obtained for them. Section
\ref{sec:testing_protocols} closes the chapter with the four measurement
protocols, latency and memory, rectification tolerance, real-camera
agreement, and the dual-axis evaluation convention, each stated once
here so that reported numbers need not re-derive the procedure behind
them.

\FloatBarrier
\section{Experimental Setup and Test Bench}
\label{sec:exp_setup_bench}
The test bench spans four compute platforms and one piece of physical
sensing hardware, chosen so that each stage of the work runs on the
platform best suited to it. The production model is implemented in
PyTorch 2.11 with CUDA 12.8 on Python 3.10, and training runs on a
single cloud GPU with 80 gigabytes of memory. Full runs are expensive
enough that they are reserved for the final configuration; the many
controlled ablations behind that configuration, and the cross-domain
evaluation that follows every training run, instead run on a smaller
cloud GPU, chosen for its low cost so that a controlled comparison after
every design change stays practical. Local latency benchmarking runs on
a laptop-class GPU with roughly 4 gigabytes of memory, which represents
the low end of the memory budget the network must fit inside rather than
a generous development machine. The intended deployment target is a
Jetson Orin Nano class of edge accelerator, on which the optimized model
was measured directly, and Figure~\ref{fig:jetson} shows the development
board used for that measurement.

\begin{figure}[!htbp]
\centering
\includegraphics[width=0.6\textwidth]{fig_7_jetson}
\caption{Jetson Orin Nano edge accelerator, deployment target for inference.}
\label{fig:jetson}
\end{figure}

On the sensing side, the test bench includes a calibrated low-cost
binocular rig, a Waveshare AR0144 stereo USB camera
\cite{waveshare_ar0144} described fully in
Section~\ref{sec:calibration}, which supplies every real-camera pair
used in this thesis. Together, the four
compute platforms and the rig give a training platform sized for the
full dataset, an inexpensive platform for iterating on the design, a
worst-case memory reference for the deployment budget, the intended
deployment hardware itself, and a physical source of real, imperfectly
rectified data that no synthetic benchmark can substitute for.

\FloatBarrier
\section{Calibration and Rectification}
\label{sec:calibration}
Real-camera pairs come from a calibrated binocular rig, a Waveshare
AR0144 stereo USB camera whose dual global-shutter sensors are driven
synchronously so both views land in the same frame
\cite{waveshare_ar0144}. Figure \ref{fig:ar0144} shows the module. The
rig streams a single 2560 by 720 side-by-side image, which is split
into a rectified 1280 by 720 pair per frame; Figure
\ref{fig:rig_schematic} sketches the stream and the calibrated
geometry, and Table \ref{tab:rig_specs} lists the manufacturer
specifications.

\begin{figure}[!htbp]
\centering
\includegraphics[width=0.72\textwidth]{fig_5_1b_ar0144_baseline}
\caption[Waveshare AR0144 stereo camera module with its 52 mm baseline.]{Waveshare AR0144 stereo camera module with its 52 mm baseline (product image: Waveshare \cite{waveshare_ar0144}).}
\label{fig:ar0144}
\end{figure}

\begin{table}[!htbp]
\centering
\caption[Manufacturer specifications of the AR0144 stereo camera.]{Manufacturer specifications of the AR0144 stereo camera \cite{waveshare_ar0144}.}
\label{tab:rig_specs}
\small
\begin{tabular}{ll}
\toprule
Item & Value\\
\midrule
Sensors & 2 $\times$ onsemi AR0144, 1/4 inch, global shutter\\
Stereo output & $2560 \times 720$ side by side ($1280 \times 720$ per view), 30 FPS\\
Baseline & 52 mm\\
Lens & 2.88 mm EFL, F/2.2, fixed focus, distortion $< 0.2\%$\\
Field of view & $74^{\circ}$ diagonal, $65^{\circ}$ horizontal, $43^{\circ}$ vertical\\
Interface & USB 2.0 Type-C, 5 V\\
Board size & $66.00 \times 30.00 \times 17.72$ mm\\
\bottomrule
\end{tabular}
\end{table}

\begin{figure}[!htbp]
\centering
\includegraphics[width=\textwidth]{fig_5_2_rig_preview}
\caption{Stereo rig capture and rectification geometry for metric depth.}
\label{fig:rig_schematic}
\end{figure}

Calibration recovers a focal length of
approximately 1005 pixels at the native 1280-pixel width and a baseline
of 52 millimeters between the two camera centers; both values scale to
the network's inference resolution by the width ratio, so the disparity
the network predicts converts to metric depth without any additional
tuning. The calibrated values agree with the manufacturer data of
Table \ref{tab:rig_specs}: the published baseline is exactly 52
millimeters, and the 65 degree horizontal field of view over the
1280-pixel width implies a focal length close to the calibrated 1005
pixels. The rig's factory rectification is imperfect and leaves a
residual vertical misalignment between the two views, on the order of a
fraction of a pixel to about a pixel, rather than the perfect
epipolar alignment synthetic training data provides. This gap between
laboratory-perfect and factory rectification is exactly the condition
the second research objective targets, and Section
\ref{sec:testing_protocols} defines the two protocols, one controlled
and synthetic, one on the physical rig, that measure how well the
network tolerates it.

\FloatBarrier
\section{Datasets}
\label{sec:exp_datasets}
This thesis draws on three data sources, summarized in Table
\ref{tab:exp_datasets}: a synthetic corpus with dense ground truth for
training and in-domain testing, a real benchmark the network never
trains on for zero-shot cross-domain testing, and the project's own
physical rig for real-camera qualification. This section describes what
each source contains.

\begin{table}[!htbp]
\centering
\caption{Data sources with pair counts, resolution, and ground truth.}
\label{tab:exp_datasets}
\small
\begin{tabular}{lllll}
\toprule
Source & Role & Pairs & Resolution & Ground truth\\
\midrule
SceneFlow finalpass & train / in-domain test & 35{,}454 / 4{,}370 & $960{\times}540$ & dense, rendered\\
Middlebury 2014 & zero-shot cross-domain & 23 (perfect set) & up to ${\sim}3000{\times}2000$ & structured light\\
Stereo rig capture & real-camera qualification & 997 & $1280{\times}720$ & none (teacher reference)\\
\bottomrule
\end{tabular}
\end{table}

\subsection{Synthetic Training and Test Corpus: SceneFlow}
SceneFlow finalpass \cite{mayer2016sceneflow} is the synthetic corpus the
network trains on. Its canonical split combines FlyingThings3D,
Monkaa, and Driving into 35,454 training pairs, each rendered at
$960\times 540$ with dense, pixel-exact disparity ground truth; only the
FlyingThings3D subset provides an official test split, 4,370 pairs, held
out and never seen during training. Chapter~\ref{ch:implementation}
describes the crop, augmentation, and optimization protocol applied on
top of this corpus; this chapter records the corpus itself, since its
size and density set the ceiling on how well the network can learn
dense correspondence before it is ever asked to generalize. Rendered
ground truth of this kind is something no physical sensor can supply at
this scale, which is precisely why the field trains on it despite the
domain gap it introduces relative to real imagery.

\subsection{Zero-Shot Cross-Domain Benchmark: Middlebury 2014}
Middlebury 2014 \cite{scharstein2014middlebury} supplies the cross-domain
test. Its 23 perfect-rectification scenes, real indoor imagery captured
with structured-light ground truth at up to roughly $3000\times 2000$
pixels, carry public ground truth and are the scenes this thesis
evaluates on; the benchmark's remaining 10 scenes withhold their ground
truth on a submission server and are not used here. The network never
trains on any Middlebury imagery, so every result against this
benchmark is a genuine zero-shot measurement.
Its disparities extend to several hundred pixels and its scenes include
thin structures and repeated texture that the synthetic training
distribution under-represents, which is why it functions as a stress
test rather than a second in-domain benchmark.

\subsection{Physical Stereo Rig Captures}
The third source is the project's own rig, described in Section
\ref{sec:calibration}. Unlike the two benchmarks above, it carries no
dense ground truth of its own; its role is to confirm that behaviour
measured on synthetic and public data survives contact with a real,
imperfectly rectified camera. Section~\ref{sec:data_collection} details
how these pairs were captured and how a reference disparity was
obtained for them despite the absence of ground truth.

\FloatBarrier
\section{Data Collection}
\label{sec:data_collection}
Real pairs were captured directly from the rig's live 2560 by 720
side-by-side stream, described in Section~\ref{sec:calibration}, and
split in software into the rectified 1280 by 720 left and right frames
that the network consumes at its inference resolution. No frame
selection criterion beyond basic scene variety was applied: the capture
set spans ordinary indoor scenes at the range the rig was calibrated
for, since the purpose of these pairs is to qualify the network on
unremarkable real data rather than to stress a particular scenario.
This yielded 997 rectified pairs, none of which contributed to training
or fine-tuning at any stage; the network that produces disparity on
these pairs is exactly the SceneFlow-trained checkpoint of Chapter
\ref{ch:implementation}.

Because the rig supplies no dense ground truth, the real-camera
agreement protocol of Section~\ref{sec:testing_protocols} instead
compares the network's prediction on these 997 pairs against a
pseudo-ground-truth reference produced by a frozen FoundationStereo
teacher run on the same pairs. The teacher is a large, accurate,
foundation-model-based network that is never trained or fine-tuned on
this project's data; it is applied once, off-line, to generate a
disparity reference for each of the 997 pairs, and that reference is
then treated as a strong stand-in for ground truth in the agreement
comparison. Because both the compact network and the teacher observe
the same imperfectly rectified input, a close match between them
measures agreement under real conditions rather than absolute error
against a laboratory-perfect reference, and Section
\ref{sec:testing_protocols} states this distinction explicitly wherever
the resulting numbers are used.

\FloatBarrier
\section{Experimental Procedure and Evaluation Metrics}
\label{sec:testing_protocols}
Every reported number follows one of four measurement protocols,
defined here so that each has a single stated procedure rather than an
implicit one repeated across tables. Accuracy
throughout follows the standard stereo metrics of \eqref{eq:metrics}:
end-point error \cite{mayer2016sceneflow}, the bad-$t$ rates
\cite{scharstein2002taxonomy}, and the D1-all outlier rate
\cite{menze2015kitti}, each metric averaged per image and then over
images rather than pooled across every pixel in the set.

The first protocol measures latency and memory. Latency is measured at
batch one and $384\times 640$ resolution after warm-up passes; memory is
recorded as the peak inference working set, and every measurement states
its device. On the local laptop GPU the model runs in eager PyTorch at
full and half precision, without compiler assistance, so those numbers
reflect an unoptimized serving path. On the deployment target, the
Jetson Orin Nano, the optimized model is exported to an eight-bit engine
and measured directly on the device, and that on-device figure is the
one taken as evidence for the real-time claim.

The second protocol tests tolerance to imperfect rectification under
controlled conditions. The trained model is evaluated at the native axis
on a fixed 400-pair FlyingThings3D subset. The right image is shifted
vertically by a controlled offset while the disparity ground truth stays
unchanged, so the offset emulates a residual vertical misalignment
between the views without touching anything else about the pair. The
sweep runs from zero through sub-pixel offsets to several pixels, and
the standard metrics are reported at every point on the sweep.

The third protocol corroborates that synthetic sweep on physical data,
using the collection described in Section~\ref{sec:data_collection}.
The trained model is run on the 997 rectified pairs from the project's
rig, whose factory rectification leaves a residual vertical misalignment
of the kind the second protocol emulates, and its disparity is compared
against the FoundationStereo teacher's pseudo-ground-truth on the same
pairs. Because both the model and the teacher see the same misaligned
input, this protocol measures agreement with a strong reference rather
than absolute error, and it corroborates rather than replaces the
controlled sweep of the second protocol.

The fourth is the dual-axis evaluation convention. Every in-domain
evaluation is reported on the native $960\times 540$ axis, which
preserves true disparity magnitudes, with a resized-axis measurement at
$640\times 384$ logged alongside for comparability with the legacy
protocol discussed in Chapter~\ref{ch:implementation}. Any comparison
between input protocols evaluates every candidate on both axes, so a
gain on one axis cannot hide a loss on the other. Zero-shot Middlebury
2014 evaluation follows a single fixed protocol applied identically to
the trained model and to the public reference checkpoints it is compared
against: inference at $384\times 640$, ground truth scaled to the
inference width, disparities above 192 pixels masked, and per-scene
metrics macro-averaged over scenes.

Each protocol runs as a separate scripted driver on a cloud GPU, and
every driver writes its configuration and metrics to a metadata file
alongside its outputs so that a reported number can be traced back to
the run that produced it; the specific driver names are recorded in the
appendix covering technical specifications. A cross-domain evaluation
run of this kind
costs only minutes and a small fraction of a cloud GPU-hour, which is
what keeps the after-every-training-run evaluation gate described in
Chapter~\ref{ch:design} practical rather than an occasional afterthought.

\FloatBarrier
\section{Summary}
This chapter fixed the test bench, the data sources, and the
measurement protocols behind every number in the study. The test bench
pairs a large cloud GPU for production training with a smaller,
cheaper cloud GPU for ablations, a memory-constrained laptop GPU for a
conservative latency reference, and the Jetson Orin Nano as the
deployment target, alongside a calibrated low-cost stereo rig whose
focal length, baseline, and imperfect factory rectification were all
stated so that reconstruction and robustness numbers stay traceable to
a physical measurement. Three data sources were
described: the SceneFlow synthetic corpus that trains and
in-domain-tests the network, the Middlebury 2014 benchmark that tests
it zero-shot on unseen real imagery, and 997 real pairs collected from
the project's own rig, for which a frozen FoundationStereo teacher
supplied a pseudo-ground-truth reference in place of the dense ground
truth the rig cannot provide. Four measurement protocols, latency and
memory, controlled rectification perturbation, real-rig agreement, and
dual-axis evaluation, were each fixed to a single stated procedure,
giving every reported number a documented and reproducible method
behind it rather than an implicit one.

% SOURCES: old mainmatter/ch7.tex Section "Simulation and Experimental
%   Setup" (7.6, full text: three data sources, fig:setup and fig:jetson
%   figures, scripted-driver reproducibility framing, low-cost-evaluation
%   sentence) and Section "Testing and Measurement" (7.8, full text: all
%   four measurement protocols, \label{sec:testing_protocols} kept
%   verbatim). old mainmatter/ch2.tex "Benchmarks and Datasets" subsection
%   (SceneFlow 35,454/4,370 pairs at 960x540, Middlebury 2014 23/10 scenes
%   up to ~3000x2000, both \cite keys reused). old mainmatter/ch3.tex
%   Principle-of-Operation paragraph (rig focal length ~1005 px at native
%   1280 px width, baseline 52 mm; numbers only). old appendices/app_a.tex
%   "Software and Hardware Environment" (PyTorch 2.11, CUDA 12.8, Python
%   3.10, per-platform memory figures) and "Reproduction Pointers"
%   (997-pair in-house capture set note, scripted-driver-per-experiment
%   framing).
% DROPPED: the fig:setup diagram (fig_7_setup_preview) from old 7.6 was
%   considered for this chapter but its content (three data sources plus
%   training/evaluation platforms plus recorded metrics) is now fully
%   covered in prose across Sections 5.2-5.6 of this chapter without the
%   figure, so it was left out here rather than duplicated; if a future
%   integration pass wants the original diagram reinstated, it is
%   unchanged at fig_7_setup_preview and carries the label fig:setup in
%   the old book. Vendor names for the cloud GPUs (a specific commercial
%   provider) are not stated, matching the vendor-neutral convention
%   already used throughout the old book's Chapter 7 and 8 text.
% MOVED-AWAY: SceneFlow crop size, augmentation triplet, optimizer,
%   schedule, batch size, precision, and Table tab:train_config -> owner
%   ch:implementation (old 7.2 Datasets and Training Protocol + 7.3
%   Selection of Parameters, per the plan's row for 7.2/7.3/7.4/7.7).
%   The disparity-to-depth back-projection mechanism and the network's
%   forward-pass description -> owner ch:implementation (old 7.7
%   Principle of Operation). The literature-survey dataset comparison
%   table tab:datasets (with KITTI 2012/2015 and ETH3D, not used in this
%   thesis's own experiments) stays with ch:lit's comparative-analysis
%   section; this chapter's Table \ref{tab:exp_datasets} is a new,
%   distinctly labeled table scoped to the three sources this thesis
%   actually runs on, to avoid a duplicate tab:datasets label at
%   integration. The controlled-ablation review-and-feedback gate
%   referenced in the "low-cost evaluation" sentence -> owner ch:design
%   (old 7.5 Review and Feedback).
